We begin by fixing the notation and the objects that the whole replication will use. Let
a transformer language model be composed of $L$ blocks, so that, for an input prompt $X$,
block $l$ has activations $A_l(X) \in \mathbb{R}^{t \times k }$, where $t$ is the sequence
length and $k$ the hidden dimension. As in Beaglehole et al., the feature associated with
a prompt is the last real token row of each block output, $a_l \in \mathbb{R}^k$. For each
concept we gather $n$ labelled prompts $\{X^{(i)}, y^{(i)}\}$, where $y \in \{0,1\}$
records the presence of the concept in each prompt, we run every prompt through the model,
and we stack the per-block last-token activations into the matrix
$Z_\ell \in \mathbb{R}^{n \times k}$ whose $i$-th \emph{row} is
$\left(a_\ell^{(i)}\right)^{\top}$. Note that extracting a concept vector requires no
generated text at all, only $(\text{prompt}, \text{label})$ pairs and one forward pass
each, which is what makes the data collection cheap and what the steering and monitoring
stages downstream then build on. It is on this cheapness that Beaglehole et al. rely when
they extract and steer a concept in under a minute with fewer than $500$ samples.

\subsection{Recursive Feature Machine}

The Recursive Feature Machine (RFM) of \cite{radhakrishnan2024mechanism} can be read,
informally, as a supervised and nonlinear counterpart to Principal Component Analysis
(hereafter PCA), in the sense that, where PCA returns the directions of greatest
\emph{variance} in the data, RFM returns the directions to which a fitted \emph{predictor}
is most sensitive, via the Average Gradient Outer Product (hereafter AGOP). It is not a
generalisation of PCA in the strict sense, since PCA is not a special case of the
iteration below, but the analogy is the right one for understanding what the matrix $M$ is
doing.

The predictor, for this case, is a kernel ridge regression (hereafter KRR) with a
Mahalanobis-Laplace kernel parametrized by a Positive Semidefinite (PSD) matrix $M$ with
bandwidth $\sigma$, given by

\[
  K_M(x,z) \;=\; \exp\!\left(-\frac{1}{\sigma}\sqrt{(x-z)^{\top} M (x-z)}\right)
\]

where $\sigma>0$ is the kernel bandwidth. Note that we write $\sigma$ where Beaglehole et
al. write $L$, so as to avoid a clash with the number of transformer blocks $L$.

Initialize $M_0 = I$ and, for $t = 0, \dots, T-1$:

\begin{enumerate}
  \item \textbf{Kernel ridge regression.} Fit
  \[
    \alpha_t = \left(K_{M_t}(Z,Z) + \lambda I\right)^{-1} y \;\in\; \mathbb{R}^{n},
  \]
  \[
    f_t(z) = \alpha_t^{\top} K_{M_t}(Z,z) = \sum_{i=1}^{n} \alpha_{t,i}\, K_{M_t}\!\left(a^{(i)}, z\right).
  \]

  \item \textbf{AGOP.}\footnote{We mean-centre the gradients before forming the outer product, as
pre-registered in \texttt{plans/GUPPYLM\_GUIDE.md} §1a and as implemented in
\texttt{rfm\_guppy/rfm.py}. Without centring, the top eigenvector of $M_{t+1}$ is pulled
toward the mean gradient direction rather than the direction of greatest gradient
variation.} Recompute $M$ as the average outer product of the predictor's
  input-gradients over the training points:
  \[
    M_{t+1} = \frac{1}{n}\sum_{i=1}^{n}
      \left(g^{(i)} - \bar{g}\right)\left(g^{(i)} - \bar{g}\right)^{\top},
  \]
  \[
    g^{(i)} = \nabla_z f_t\!\left(a^{(i)}\right),
    \quad \bar{g} = \frac{1}{n}\sum_{j=1}^{n} g^{(j)} .
  \]
\end{enumerate}

After $T$ iterations, the concept vector $v$ is taken as the top eigenvector of $M_T$,
with $\|v\|_2 = 1$. It is then oriented by computing the Pearson correlation $\rho$
between the projections $\{\langle a^{(i)}, v\rangle\}$ and the labels $\{y^{(i)}\}$ and
by replacing $v$ with $\operatorname{sign}(\rho)\, v$, so that a positive projection means
``more concept''. The gradient that drives the AGOP, writing
$d_i(z) = \sqrt{(z - a^{(i)})^{\top} M (z - a^{(i)})}$, is then

\[
  \nabla_z f_t(z) \;=\; -\frac{1}{\sigma}\sum_{i=1}^{n}
    \alpha_{t,i}\, K_{M_t}\!\left(a^{(i)}, z\right)\,
    \frac{1}{d_i(z)}\; M\left(z - a^{(i)}\right),
\]

with the convention that terms for which $d_i(z) = 0$ are dropped. Typical
hyperparameters are $\lambda = 10^{-3}$, $\sigma \in \{1, 10, 100\}$ and $T \leq 10$.

We prefer RFM to a plain logistic regression because the AGOP reports the directions to
which a \emph{nonlinear} predictor is sensitive, so that it can in principle surface
concept structure that a single linear separator misses. Beaglehole et al. advance that
claim under two conditions, namely abstract concepts and large models, and it is worth
keeping the conditions attached, since they decide the outcome of
§\ref{sec:baselines-anchor} below, where on the three easy and in-domain concepts
available here logistic regression matches and slightly exceeds RFM.

\subsection{Steering}

Given the concept vectors $\{v_l\}$ obtained per block, we steer during the forward pass
by adding $\varepsilon v_l$ to every token row of the output of block $l$, that is,

\[
  \tilde{A}_{\ell,i}(X) = A_{\ell,i}(X) + \varepsilon\, v_\ell
  \qquad \text{for every token position } i,
\]

where $\varepsilon$ is the control coefficient and the concept vector is taken with
$\|v_\ell\|_2 = 1$. Note that the sign of $\varepsilon$ carries the direction of the
intervention, since a positive $\varepsilon$ steers towards the concept and a negative one
steers away from it. If one wishes to steer a group of concepts simultaneously, one shall
use a combination of the corresponding vectors.

\subsection{Monitoring}

For monitoring rather than steering, we keep the top $p$ eigenvectors $\{v_{l,j}\}$ per
block instead of a single one. The feature for prompt $i$ is then the set of projections
$\langle a_l^{(i)}, v_{l,j} \rangle$ over all blocks $l$ and all components $j$, which is
a vector in $\mathbb{R}^{Lp}$. On these features we train a classifier under one of two
strategies, either per-block, training one classifier per block and keeping the best of
them, or aggregate, concatenating all blocks and training a single classifier. The
evaluation is made by AUROC (Area Under the Receiver Operating Characteristic curve).

\subsection{Baselines}

RFM is compared against three cheaper extraction methods, described below. All three
operate on the mean-centred training activations, return a unit vector, and are oriented
by the same Pearson-sign rule as RFM, so that the comparison between them is a comparison
of directions and of nothing else.

\begin{itemize}
  \item \textbf{PCA.} Pair the $j$-th positive sample with the $j$-th negative sample,
  $j = 1, \dots, m$ with $m = \min(n_+, n_-)$, and form the difference matrix $D$ with
  rows $D_j = a^{(\pi^+(j))} - a^{(\pi^-(j))}$. Mean-centre it, $\tilde{D} = D - \bar{D}$,
  so that $v$ is the top right singular vector of $\tilde{D}$, equivalently the top
  eigenvector of $\tilde{D}^{\top}\tilde{D}$.

  \item \textbf{Difference-in-means.}
  $v = \dfrac{\mu_1 - \mu_0}{\left\|\mu_1 - \mu_0\right\|_2}$, where
  $\mu_c = \operatorname{mean}(S_c)$ is the mean activation of class $c$.

  \item \textbf{Logistic regression.} $v = w / \|w\|_2$, for the fitted coefficient
  vector $w$, with the inverse-regularisation strength $C$ swept on the validation split.
\end{itemize}
